%%=============================================================================
%% Validatie en resultaten
%%=============================================================================

\chapter{\IfLanguageName{dutch}{Validatie en resultaten}{Validation and results}}%
\label{ch:validatie}

In dit hoofdstuk worden de resultaten besproken van de validatie van de PoC uit Hoofdstuk~\ref{ch:poc}
en, indirect, van het framework als geheel. De validatie volgt de aanpak die in \emph{Component~4} van het
framework (\autoref{sec:component-validatie}) werd vastgelegd: een Parallel Run met meerdere testscenario's,
aangevuld met een expert-review.

\section{Opzet van de Parallel Run}
\label{sec:val-opzet}

De Parallel Run werd uitgevoerd in een gedeelde testomgeving op het mainframe waarin zowel de
oorspronkelijke CA~Gen-module \texttt{TPPBAMT} als de nieuwe PL/I-implementatie aanwezig zijn. Beide
implementaties werden achtereenvolgens gestart op exact dezelfde DB2-snapshot en met identieke
inputparameters. De gegenereerde uitvoerbestanden worden bewaard, zodat ze achteraf reproduceerbaar konden 
worden vergeleken.

\section{Testdatasets}
\label{sec:val-datasets}

Conform \autoref{subsec:test-datasets} werden drie scenario's voorzien:
\begin{description}
    \item[Scenario~1 -- Standaard] Een dagelijkse productieaanvraag met een representatief volume aan
        planbare materialen voor \'e\'en oven.
    \item[Scenario~2 -- Edge-case] Een aanvraag die de maximumkardinaliteit van de Group View benadert,
        gecombineerd met enkele uitzonderlijke materiaalkwaliteiten en sequentievoorwaarden.
    \item[Scenario~3 -- Foutscenario] Een aanvraag met onvolledige of inconsistente input, bedoeld om de
        foutpaden en exitstate-mapping te valideren.
\end{description}

\section{Resultaten}
\label{sec:val-resultaten}


\section{Expert-review}
\label{sec:val-expert}



\section{Reflectie op het framework}
\label{sec:val-reflectie-framework}

